% Draft for author review. No revised outcome claims belong in this section.
\section{Introduction}
\label{sec:introduction}

A voltage surrogate can approve an operating point quickly. The relevant question is whether its decision remains useful after the feeder changes. A learned predictor trained on one network need not receive the information that distinguishes two switch configurations. When that information is absent, repeated predictions on the same loads, photovoltaic injections and slack-voltage setting can remain confident even though the physical voltage profile changes.

Exact power-flow fallback does not settle which points should be checked. Escalating every case removes the computational purpose of the surrogate. Escalating too few can leave unsafe approvals, while rejecting uncertain cases can discard feasible operation. These costs must be measured together. A physics-assisted neural route therefore requires comparison with the same physics model used directly, with exact fallback, to identify what the neural component contributes.

We study this comparison on a synthetic, balanced 33-bus radial feeder. The stale neural ensemble sees operating features but no current topology. LinDistFlow receives the supplied network model. A separate residual predictor also receives the current LinDistFlow voltage extrema and is trained across development topologies. This information difference is intentional and explicit. The study asks when external physics helps a stale predictor and whether a learned correction improves on physics alone. Models with and without topology information are different observers.

The evaluation pairs operating points across every connected radial single branch exchange in the feeder inventory. It separates development exchanges from held-out exchanges, fixes thresholds using independent calibration data and perturbs the verifier while keeping the physical reference unchanged. It preserves unresolved reference solves. These choices distinguish ranking quality, frozen-route transfer and dependence on accurate external information.

\subsection{Related work and scope}

Reject-option prediction formalizes the tradeoff between error and refusal~\cite{chow1970}. Learning to defer extends this setting to a downstream expert~\cite{mozannar2020}, while cascade analysis explains why the first model's confidence need not identify when another model helps~\cite{jitkrittum2023}. Deep ensembles motivate a spread diagnostic~\cite{lakshminarayanan2017}; the present member-margin spread is not a reproduction of every component of that method. Calibrate-Then-Delegate studies risk and budget guarantees for cascades~\cite{pona2026}. Our local probes do not implement those guarantees. We instead include a separate split-conformal margin interval and state its exchangeability requirement~\cite{angelopoulos2023}.

Combining a complex learned component with simpler verification has an established architectural history~\cite{sha2001}. Power-system research also addresses feasibility directly. DeepOPF+ calibrates constraints during training~\cite{zhao2020}, and Venzke et al. study worst-case guarantees for neural optimal power flow~\cite{venzke2020}. Parker constructs adversarial AC-power-flow inputs that distinguish surrogate feasibility from physical feasibility~\cite{parker2026}. These results motivate checking decisions against the governing equations rather than interpreting predictive confidence as a feasibility certificate.

Exact fallback is already part of this literature. Shamseldein combines learned AC feasibility, conformal prediction and an exact Newton--Raphson solver~\cite{shamseldein2026}. Nawaz et al. describe physics-informed graph surrogates with residual-triggered Newton--Raphson fallback and topology perturbation tests~\cite{nawaz2026}. Our contribution is therefore a controlled comparison of routing choices and their information requirements, not the first use of physics or exact fallback. Its scope is one feeder family, specified synthetic operating distributions and enumerated branch exchanges. Any advantage must survive comparison with physics-only routes, false-rejection accounting and measured computation.
